%% Interstitial: Paper 1 -> Paper 2
%% Narrative: the probability is determined by its marginals;
%% now we ask what predictive structure it generates.

Chapter~\ref{chap:paper-one} established that a probability measure
$P$ on $(\Omega, \sigma(\mathcal{Q}))$ is completely determined by its
query marginals.  What the observable interface shows is all there is.

This raises the next question: what does the observable interface
\emph{predict}?  Given the current query $Q$ and some future
observable $F$, the probability $P$ determines a conditional
distribution of $F$ given $Q$.  But this conditional distribution
depends on the full measure $P$, which in turn depends on the entire
query system.  Can the predictive content of $Q$ be isolated ---
expressed as a structure intrinsic to $Q$ and $P$, without reference
to all other queries?

\medskip

Chapter~\ref{chap:paper-two} answers yes.  The \emph{predictive
kernel} $\Pi_Q$ captures exactly the information in $Q$ that is
relevant for predicting $F$: it is the conditional distribution of
$F$ given $Q$, factored through the \emph{minimal predictive query}
$Q_* = \varphi_Q \circ Q$ --- the coarsest function of $Q$ that
retains all predictive content.

The central result is the \emph{predictive factorization theorem}:
every bounded measurable function $g$ of $F$ satisfies
\[
\mathbb{E}[g(F) \mid \sigma(Q)] = (\text{predictiveOp}\, g) \circ Q_*.
\]
Prediction factors through the minimal predictive state.  The query
system does not need to be refined further to predict $F$; the
sufficient statistic $Q_*$ already carries everything.

\medskip

This is the first result in the program that is genuinely dynamical
in character.  The query $Q$ represents the observer's current
measurement; $F$ represents a future event; $Q_*$ is the minimal
structure needed to bridge them.  The passage from observable
compatibility to predictive sufficiency is the passage from static
coexistence of queries to the temporal structure of prediction.

\medskip

The passage from Chapter~\ref{chap:paper-one} to
Chapter~\ref{chap:paper-two} is the passage from \emph{determination}
to \emph{prediction}: from knowing that $P$ is fixed by the query
system to extracting, for each query, the minimal predictive structure
it contains.
